\section{Conclusion}

For an assistant limited to commune-level procedures, the hard part of real citizen questions is deciding whether to answer at all. A
system that looked strong on team-written questions answered 114 of the 248 real questions it should have refused. A small, transparent
rule layer nearly doubled the end-to-end correct answers on a locked test set and kept its gain in refusals on a second one; neither a
reranker nor a local 4B LLM reader significantly improved balanced or end-to-end accuracy, nor did two gates learned \emph{post hoc} in place
of the hand-written one, and letting the model decide was no better than always refusing. A second round of rules that looked better on
development tied the first on new questions, and answers remain wrong
more often than right on random traffic. Future work: confidence-based answering with human confirmation, procedure names proposed by
the model while rules decide, a third locked set, and independent human annotation.
